\section{VOA reticular, involución canónica y rama Moonshine de orden dos}
\label{sec:voa-moonshine-orden-dos-rev11}
\index{VOA!retículo de Leech}
\index{Moonshine!clase 2A}

La rama de incidencia construye primero el retículo de Leech
\(\Lambda_{24}\). A partir de ese dato, la teoría clásica de álgebras de
operadores de vértice proporciona una prolongación canónica. Si
\(\mathfrak h=\Lambda_{24}\otimes_{\mathbb Z}\mathbb C\) y \(M(1)\) es
el módulo de Fock de su álgebra de Heisenberg, el VOA reticular es
\begin{equation}
 V_{\Lambda}=M(1)\otimes\mathbb C[\Lambda_{24}],
 \label{eq:voa-reticular-leech-rev11}
\end{equation}
con carga central \(24\). Para \(\alpha\in\Lambda_{24}\), sus operadores
de vértice se escriben
\[
 Y(e^\alpha,z)=E^-(-\alpha,z)E^+(-\alpha,z)e^\alpha z^{\alpha_0}.
\]
La ausencia de raíces en \(\Lambda_{24}\) implica la anulación del sector
de peso uno en el orbifold de Moonshine.

La involución
\[
 \theta:\Lambda_{24}\longrightarrow\Lambda_{24},
 \qquad \theta(\lambda)=-\lambda,
\]
es intrínseca y se prolonga a \(V_{\Lambda}\). Su traza graduada se obtiene
por conteo directo de los modos osciladores.

\begin{theorem}[Traza de la involución de negación]
\label{thm:traza-theta-leech-rev11}
\begin{equation}
 t_2(\tau)
 :=\operatorname{Tr}_{V_\Lambda}
 \bigl(\theta q^{L_0-1}\bigr)
 =q^{-1}\prod_{n\geq1}(1+q^n)^{-24}.
 \label{eq:t2-producto-leech-rev11}
\end{equation}
\end{theorem}
\begin{proof}
La involución empareja \(e^\lambda\) con \(e^{-\lambda}\); como el retículo
es libre de torsión, sólo \(\lambda=0\) contribuye a la traza del sector
reticular. Cada uno de los veinticuatro osciladores de energía \(n\) aporta
\(\sum_{k\geq0}(-1)^kq^{nk}=(1+q^n)^{-1}\). El desplazamiento del vacío
produce el factor \(q^{-1}\).
\end{proof}

La completación de Fricke asociada a esta traza es
\begin{equation}
 T_{2A}(\tau)=t_2(\tau)+24+\frac{2^{12}}{t_2(\tau)}
 =q^{-1}+4372q+96256q^2+1240002q^3+\cdots.
 \label{eq:T2A-hmt-rev11}
\end{equation}
El término \(24\) cancela el término constante de \(t_2\); el factor
\(2^{12}\) es la degeneración del sector torcido de veinticuatro bosones.
La misma serie de borde determina el invariante modular mediante
\begin{equation}
 \boxed{j(\tau)=\frac{(t_2(\tau)+256)^3}{t_2(\tau)^2}}.
 \label{eq:j-desde-t2-rev11}
\end{equation}
Las ecuaciones \eqref{eq:t2-producto-leech-rev11}--
\eqref{eq:j-desde-t2-rev11} construyen la rama canónica de orden dos sin
parámetros continuos.

El orbifold de Frenkel--Lepowsky--Meurman se define por
\[
 V^\natural=(V_\Lambda)^+\oplus(V_\Lambda^T)^+.
\]
Es un VOA holomorfo de carga central \(24\), satisface
\((V^\natural)_1=0\), posee carácter
\[
 \operatorname{ch}_{V^\natural}(\tau)=j(\tau)-744
\]
y su grupo de automorfismos es el grupo Monstruo
\cite{FLM1988,ConwayNorton1979,Borcherds1992}. En esta cadena HMT fija el
retículo de Leech y la involución \(\theta\); los teoremas clásicos de VOA y
FLM identifican la prolongación resultante con \(V^\natural\) y su grupo de
automorfismos. La afirmación no requiere una acción del Monstruo sobre las
cifras: la evaluación arquimediana y la rama excepcional descienden de una
preimagen HMT común y conservan invariantes diferentes.

La función vecina
\(T_{3A}=t_3+12+729/t_3\) pertenece a la teoría clásica de Hauptmoduln de
orden tres. La realización HMT activa en el corredor trítico es la clase
\(3B\); una identificación con \(3A\) requiere declarar el levantamiento específico
que cambie de clase.
